\documentclass[11pt,a4paper]{article}

\usepackage[T1]{fontenc}
\usepackage[utf8]{inputenc}
\usepackage[margin=1in]{geometry}
\usepackage{amsmath,amssymb}
\usepackage{graphicx}
\usepackage{booktabs}
\usepackage{array}
\usepackage{caption}
\usepackage{subcaption}
\usepackage{hyperref}
\usepackage{xcolor}
\usepackage{fancyvrb}
\usepackage{enumitem}
\usepackage{float}
\usepackage{parskip}

\hypersetup{
    colorlinks=true,
    linkcolor=blue!60!black,
    urlcolor=blue!60!black,
    citecolor=blue!60!black
}

\DefineVerbatimEnvironment{codeblock}{Verbatim}{
    fontsize=\small,
    frame=single,
    framesep=3pt,
    rulecolor=\color{gray!60}
}

\title{\textbf{INDENG 231 Course Project 1}\\[0.4em]
\large Building a Backtesting Simulation System for Trading Strategies}
\author{Anna Chen}
\date{May 2026}

\begin{document}
\maketitle
\tableofcontents
\newpage

% =============================================================================
\section{System Design and Assumptions}
% =============================================================================

\subsection{What I Built}

The goal of this project was to build a backtesting system from scratch --- something
that lets you try out a trading strategy on five years of historical Nasdaq-100 data
and immediately see how it would have performed.
Rather than cramming everything into one giant script, I split the code into focused
modules: one for loading data, one for the engine that runs the simulation, one for
computing performance metrics, and a folder for all the trading strategies.
This way, if you want to test a brand-new strategy, you only need to write the strategy
itself. Everything else --- the simulation loop, the trade execution, the output charts ---
just works without any changes.

\subsection{Ground Rules (Assumptions)}

Before diving into the results, here are the constraints built into the system:

\begin{itemize}[leftmargin=*]
    \item \textbf{Trades happen at the closing price.}
          Every buy or sell is executed at that day's closing price.
          There's no simulated slippage or intraday movement.

    \item \textbf{No peeking at the future.}
          When the strategy decides what to hold on day $t$, it only sees prices
          up to and including day $t$. Future prices are completely hidden.

    \item \textbf{No short selling.}
          All portfolio weights are clamped to be non-negative.

    \item \textbf{No leverage.}
          Weights are capped so they always sum to at most 1.
          Anything leftover stays in cash.

    \item \textbf{Cash earns nothing.}
          Uninvested money just sits there at a 0\% interest rate.

    \item \textbf{Transaction costs are optional.}
          The engine supports a flat-rate cost on traded notional,
          but all experiments in this report use zero transaction cost
          to keep things clean and comparable.

    \item \textbf{Stocks with missing prices are skipped.}
          If a stock has no closing price on a given day, its existing shares
          are carried over and no trade is made.
\end{itemize}

\subsection{How the Code is Organised}

Here's a quick map of the project:

\begin{codeblock}
Final Project/
|-- run_backtest.py               # Run this to reproduce all results
|-- requirements.txt
|-- nasdaq100_daily_5y.csv        # Raw data
|-- results/                      # All output charts and tables go here
\-- backtester/
    |-- data.py                   # Loads and reshapes the CSV
    |-- engine.py                 # The simulation loop
    |-- metrics.py                # Sharpe, drawdown, win rate, etc.
    \-- strategies/
        |-- base.py               # The Strategy interface (abstract class)
        |-- single_stock.py       # Momentum, Mean Reversion, RSI, Bollinger, MACD
        |-- portfolio.py          # SMA crossover, Trailing return (+weighting variants)
        \-- advanced.py           # The two new strategies from Deliverable 5
\end{codeblock}

\paragraph{How the pieces fit together.}
Figure~\ref{fig:arch} shows the flow of data through the system.
The \texttt{DataLoader} reads the CSV file and turns it into a clean price table
(rows = dates, columns = tickers).
The \texttt{BacktestEngine} then steps through each trading day one at a time.
On each day, it hands the strategy a slice of the price history up to that date,
gets back a dictionary of target weights, and executes the implied trades at the
closing price.
After the full simulation, performance metrics are computed and everything gets
saved to the \texttt{results/} folder.

\begin{figure}[H]
\centering
\begin{minipage}{0.82\textwidth}
\small
\begin{codeblock}
CSV file
   |
   v
DataLoader  --(prices table)-->  BacktestEngine
                                       |
                           for each trading day t:
                             prices[:t] --> Strategy.generate_weights()
                                       |
                                  target weights
                                       |
                              execute trades at closing price
                              record portfolio NAV
                                       |
                                       v
                              NAV series --> metrics.py --> CSV / PNG
\end{codeblock}
\end{minipage}
\caption{Data flow through the backtesting system.}
\label{fig:arch}
\end{figure}

\subsection{Adding a New Strategy}

The key design decision was to define a single, simple interface that every strategy
must follow.
To add a new strategy, you just create a Python class that inherits from
\texttt{Strategy} and implement one method:

\begin{codeblock}
class Strategy(ABC):
    @abstractmethod
    def generate_weights(self, prices: pd.DataFrame, current_date) -> dict:
        """Return {ticker: weight}, weights in [0,1], sum <= 1."""
        ...
\end{codeblock}

That's it. The engine, the metrics, and the reporting code don't need to know
anything about the strategy's internals.

% =============================================================================
\section{How to Read the Results}
% =============================================================================

Every strategy gets evaluated on the same six metrics:

\begin{itemize}[leftmargin=*]
    \item \textbf{Cumulative Return}: how much money you'd have made in total
          over the full five years, expressed as a percentage gain:
          $\text{NAV}_T / \text{NAV}_0 - 1$.

    \item \textbf{Annualised Return}: the equivalent average yearly return,
          $(1 + R_{\text{cum}})^{252/T} - 1$, where $T$ is the number of
          trading days in the backtest.

    \item \textbf{Annualised Volatility}: how much the daily returns bounced
          around, scaled up to a yearly figure ($\text{daily std} \times \sqrt{252}$).

    \item \textbf{Sharpe Ratio}: the main summary metric --- return per unit of
          risk, $(\mu_{\text{ann}} - r_f) / \sigma_{\text{ann}}$ with $r_f = 0$.
          Higher is better.

    \item \textbf{Maximum Drawdown}: the worst peak-to-trough drop in portfolio
          value during the backtest,
          $\min_t (\text{NAV}_t / \max_{s \le t} \text{NAV}_s - 1)$.
          Think of this as ``the worst you would have felt holding this strategy.''

    \item \textbf{Win Rate}: the fraction of days where the portfolio went up.
          Useful context, but not the whole story --- a strategy that wins
          only 35\% of days can still be very profitable if the winning days
          are much bigger than the losing ones.
\end{itemize}

The Sharpe ratio is probably the most useful single number because it captures
both return \emph{and} risk.
Maximum drawdown is the most gut-check relevant number --- most real investors
would stop following a strategy long before it hits its worst drawdown.

% =============================================================================
\section{Deliverable 3: Testing Strategies on a Single Stock (NVDA)}
% =============================================================================

\subsection{Why NVDA?}

I picked NVIDIA (NVDA) as the test stock because it's one of the most interesting
stories of the past five years: a chip company that became the poster child for the
AI boom and saw its stock price increase by over 10x from 2021 to 2026.
That kind of sustained, dramatic trend makes it a great stress test for seeing
which strategies can actually ride a winner versus which ones get left behind.

\subsection{The Five Strategies}

\paragraph{1. Momentum (two versions).}
The simplest trend-following idea: if the stock is above its moving average,
buy it; otherwise hold cash.
I tested two window lengths --- a short 20-day MA and a longer 50-day MA ---
to see how sensitive the results are to that choice.

\paragraph{2. Mean Reversion.}
The opposite bet: if the stock has dropped a lot recently (specifically, more than
one standard deviation below its 20-day average), buy it expecting a bounce.
Exit back to cash once the price recovers above the average.

\paragraph{3. RSI (Relative Strength Index).}
Another mean-reversion approach, this time using RSI instead of z-scores.
The idea: buy when the 14-day RSI drops below 30 (the stock is ``oversold''),
and sell when it climbs above 70 (``overbought'').

\paragraph{4. Bollinger Bands.}
Buy when the price falls below the lower Bollinger Band (the 20-day average
minus two standard deviations), which signals an unusually large dip.
Exit when the price recovers back to the middle band.

\paragraph{5. MACD Crossover.}
A classic momentum-meets-trend-detection tool.
We use the standard 12/26/9 settings: go long when the MACD line crosses above
its signal line, and exit when it crosses back below.

\subsection{Results}

\begin{table}[H]
\centering
\caption{Single-stock backtesting results for NVDA (Apr 2021 -- Apr 2026).}
\label{tab:d3}
\begin{tabular}{lrrrrrr}
\toprule
Strategy & Cum.\ Ret. & Ann.\ Ret. & Ann.\ Vol. & Sharpe & Max DD & Win Rate \\
\midrule
Momentum (MA-20)              & 596.82\% & 47.67\% & 36.88\% & 1.2927 & $-$52.38\% & 32.38\% \\
Momentum (MA-50)              & 512.56\% & 43.90\% & 38.41\% & 1.1430 & $-$32.13\% & 34.05\% \\
Mean Reversion (z=1)          & 120.58\% & 17.22\% & 31.87\% & 0.5401 & $-$45.96\% & 16.83\% \\
RSI (14, 30/70)               & 255.33\% & 28.99\% & 31.25\% & 0.9279 & $-$35.65\% & 17.54\% \\
Bollinger Band (20, 2$\sigma$) & 216.17\% & 26.00\% & 26.13\% & 0.9951 & $-$39.98\% &  9.89\% \\
MACD (12, 26, 9)              & 256.72\% & 29.09\% & 35.90\% & 0.8104 & $-$52.19\% & 26.95\% \\
\bottomrule
\end{tabular}
\end{table}

\begin{figure}[H]
\centering
\includegraphics[width=0.92\textwidth]{results/d3_single_stock_nav.png}
\caption{Portfolio value over time for each single-stock strategy on NVDA
         (normalised to 1 at the start).}
\label{fig:d3}
\end{figure}

\subsection{What the Numbers Tell Us}

\textbf{Momentum is the clear winner --- by a lot.}
The 20-day momentum strategy returned 597\% over five years with a Sharpe of 1.29.
This makes sense: NVDA spent most of this period in a strong uptrend, so a strategy
that simply asks ``is the stock above its recent average?'' and stays long when the
answer is yes captured most of that move.

\textbf{Mean reversion struggled badly.}
A stock climbing 10x doesn't ``revert.'' The mean-reversion strategy kept waiting
for big dips to buy, spending a lot of time in cash, and only managed 120\% total.
The lesson here is that mean reversion only makes sense for stocks that actually
bounce around a stable price level --- not for NVDA.

\textbf{RSI and Bollinger Bands landed in the middle.}
Both are essentially ways of buying dips, so they do capture some upside during
pullbacks, but they miss the big runs. The Bollinger Band strategy has the lowest
win rate of all (under 10\% of days are profitable) because it only enters on extreme
dips --- but when it does, the trades tend to pay off well. Its Sharpe (0.99) is
actually pretty competitive.

\textbf{MACD was decent but noisy.}
The MACD signal is slow to react because it uses a 26-day lagging EMA, so it
often enters trends late. Returns are similar to RSI but with higher volatility,
which drags the Sharpe down to 0.81.

% =============================================================================
\section{Deliverable 4: Building a Full Portfolio}
% =============================================================================

\subsection{From One Stock to 101}

Single-stock backtesting tells you how a strategy performs on one name,
but the more interesting (and realistic) question is: what happens when
you apply it across the whole Nasdaq-100 and turn the signals into a portfolio?

I tested two different ways of selecting stocks and two different ways
of deciding how much to put into each one, giving four combinations in total.

\subsubsection{Stock Selection Signals}

\paragraph{SMA Crossover.}
For each of the 101 stocks, compute a 20-day and a 50-day simple moving average.
If the short-term average is above the long-term one, that stock gets included in
the portfolio. This is a classic ``trend is your friend'' filter.

\paragraph{Trailing Return Momentum.}
Look at each stock's return over the past 30 days, rank all 101 stocks,
and take the top 10.
This is a straightforward cross-sectional momentum rule: bet on what's been
working lately.

\subsubsection{Weighting Schemes}

\paragraph{Uniform weighting.}
Split the portfolio equally across all selected stocks:
\[
w_{i,t} = \frac{1}{|\mathcal{S}_t|}, \quad i \in \mathcal{S}_t.
\]

\paragraph{Risk-adjusted (inverse-volatility) weighting.}
Instead of equal weights, allocate more to stable stocks and less to volatile ones.
Specifically, each stock's weight is proportional to the inverse of its recent
30-day volatility:
\[
\tilde{w}_{i,t} = \frac{1}{\hat{\sigma}_{i,t}}, \qquad
w_{i,t} = \frac{\tilde{w}_{i,t}}{\sum_{j \in \mathcal{S}_t} \tilde{w}_{j,t}}.
\]
The idea is to prevent a handful of wild stocks from dominating the portfolio.

\subsection{Results}

\begin{table}[H]
\centering
\caption{Portfolio backtesting results across all 101 Nasdaq-100 stocks
         (Apr 2021 -- Apr 2026).}
\label{tab:d4}
\begin{tabular}{lrrrrrr}
\toprule
Strategy & Cum.\ Ret. & Ann.\ Ret. & Ann.\ Vol. & Sharpe & Max DD & Win Rate \\
\midrule
SMA Crossover + Uniform      &  72.00\% & 11.51\% & 18.98\% & 0.6061 & $-$37.71\% & 51.67\% \\
SMA Crossover + Risk-Adj.    &  52.55\% &  8.85\% & 16.13\% & 0.5485 & $-$32.92\% & 51.04\% \\
Trailing Return + Uniform    & 239.09\% & 27.79\% & 27.60\% & 1.0068 & $-$43.25\% & 52.55\% \\
Trailing Return + Risk-Adj.  & 158.27\% & 20.99\% & 25.44\% & 0.8249 & $-$44.00\% & 52.31\% \\
\bottomrule
\end{tabular}
\end{table}

\begin{figure}[H]
\centering
\includegraphics[width=0.92\textwidth]{results/d4_portfolio_nav.png}
\caption{Portfolio value over time for each portfolio-level strategy.}
\label{fig:d4}
\end{figure}

\subsection{Key Takeaways}

\textbf{Picking the right stocks matters more than how you weight them.}
Whether you use uniform or risk-adjusted weights, the trailing-return signal
comfortably beats the SMA crossover signal. SMA crossover tends to select a broad
swath of stocks (sometimes 50+ at once), which dilutes the signal. Trailing return
picks just the top 10, keeping the portfolio concentrated in whatever's actually
working.

\textbf{Risk-adjusted weighting actually hurt performance here.}
This one is a bit counterintuitive. You'd expect inverse-vol weighting to help by
reducing exposure to risky names --- and it does lower volatility (slightly).
But in this particular period, the high-volatility Nasdaq stocks (NVDA, TSLA, MSTR, etc.)
were also the ones with the best returns. By underweighting them, the risk-adjusted
portfolios left a lot of return on the table without getting a commensurate reduction
in drawdown.

\textbf{Bottom line: Trailing Return + Uniform is the best portfolio here at Sharpe 1.01.}
This becomes the harder of the two benchmarks to beat in Deliverable 5.

% =============================================================================
\section{Deliverable 5: Two New Strategies That Beat the Benchmarks}
% =============================================================================

\subsection{The Benchmarks}

\paragraph{Benchmark 1 --- SMA Crossover.}
Select stocks where the 20-day SMA is above the 50-day SMA; allocate equally.
This is a broad market-participation strategy. \textbf{Sharpe: 0.61.}

\paragraph{Benchmark 2 --- Trailing Return Momentum.}
Pick the top 10 Nasdaq-100 stocks by 30-day return; allocate equally.
This is the tougher benchmark. \textbf{Sharpe: 1.01.}

The goal is to beat \emph{both} in terms of Sharpe ratio.

\subsection{New Strategy 1: Risk-Adjusted Momentum}

\subsubsection{The Idea}

Benchmark 2 simply asks ``which stocks went up the most?''
But that's not quite the right question.
A stock that gained 40\% in a month by whipping up and down every day is
a very different proposition from one that steadily climbed 40\% with calm,
consistent daily moves.

The insight is to rank stocks not by raw return, but by their \emph{risk-adjusted} return:
return divided by volatility, which is essentially a short-term Sharpe score for each stock.
This way, volatile lottery-ticket stocks get penalised, and we end up selecting
names whose gains were more reliable and sustainable.

\subsubsection{How It Works}
\begin{enumerate}[leftmargin=*]
    \item Compute each stock's 30-day trailing return $r_i$ and
          30-day rolling volatility $\hat{\sigma}_i$.
    \item Score each stock: $s_i = r_i / \hat{\sigma}_i$.
    \item Take the top 10 by score.
    \item Split equally among them.
\end{enumerate}

\subsection{New Strategy 2: Regime-Filtered Momentum}

\subsubsection{The Idea}

One glaring weakness of pure momentum strategies is that they stay fully invested
even when the whole market is falling.
During the 2022 Nasdaq selloff, for instance, picking the ``best'' momentum stocks
still meant riding the market down with everyone else.

The fix is to add a \emph{market regime check} before doing anything.
The idea: look at the entire Nasdaq-100 universe and count how many stocks are above
their 100-day moving average.
If fewer than 40\% are (meaning most stocks are in a downtrend), the portfolio goes
entirely to cash and waits.
Only when the market looks broadly healthy does the strategy go back to picking stocks.

\subsubsection{How It Works}
\begin{enumerate}[leftmargin=*]
    \item \textbf{Check the regime.}
          Compute the 100-day SMA for every stock.
          If fewer than 40\% of stocks trade above their 100-day SMA, hold all cash
          and skip the rest.
    \item \textbf{Select stocks (only in a bull regime).}
          Use the same risk-adjusted score $s_i = r_i / \hat{\sigma}_i$ to pick
          the top 10 stocks.
    \item \textbf{Weight by inverse volatility.}
          Allocate more to the calmer stocks: $w_i \propto 1/\hat{\sigma}_i$.
\end{enumerate}

The breadth indicator (fraction of stocks above their long-run average) is a simple
but reliable signal of market-wide health.
When it flips bearish, the strategy steps aside rather than fighting the tape.

\subsection{Results}

\begin{table}[H]
\centering
\caption{New strategies vs.\ benchmarks (Apr 2021 -- Apr 2026).}
\label{tab:d5}
\begin{tabular}{lrrrrrr}
\toprule
Strategy & Cum.\ Ret. & Ann.\ Ret. & Ann.\ Vol. & Sharpe & Max DD & Win Rate \\
\midrule
Benchmark 1 --- SMA Crossover   &  72.00\% & 11.51\% & 18.98\% & 0.6061 & $-$37.71\% & 51.67\% \\
Benchmark 2 --- Trailing Return & 239.09\% & 27.79\% & 27.60\% & 1.0068 & $-$43.25\% & 52.55\% \\
\midrule
\textbf{New 1 --- Risk-Adj.\ Momentum}  & \textbf{217.26\%} & \textbf{26.09\%} & \textbf{22.49\%} & \textbf{1.1598} & $-$31.27\% & 51.83\% \\
\textbf{New 2 --- Regime-Filtered Mom.} & \textbf{142.30\%} & \textbf{19.45\%} & \textbf{16.71\%} & \textbf{1.1639} & $-$13.15\% & 40.91\% \\
\bottomrule
\end{tabular}
\end{table}

\begin{figure}[H]
\centering
\includegraphics[width=0.92\textwidth]{results/d5_benchmark_nav.png}
\caption{Portfolio value over time: new strategies vs.\ benchmarks.}
\label{fig:d5}
\end{figure}

\subsection{Did It Work?}

Yes --- both new strategies beat both benchmarks on Sharpe ratio.

\textbf{Strategy 1 (Risk-Adjusted Momentum)} gets there by keeping returns high
while cutting volatility.
It earned 217\% cumulative return (not far behind Benchmark 2's 239\%),
but with noticeably lower volatility: 22.5\% versus 27.6\%.
The result is a Sharpe of 1.16 versus 1.01.
The strategy is essentially saying: ``I still want to ride momentum, I just want
to be smarter about \emph{which} momentum stocks I pick.''

\textbf{Strategy 2 (Regime-Filtered Momentum)} takes a completely different approach.
It gives up some return (142\% vs 239\%) in exchange for dramatically lower risk.
Its maximum drawdown is only $-13.15\%$ --- compared to $-43.25\%$ for Benchmark 2.
That's a 70\% reduction in worst-case loss.
The regime filter worked exactly as intended: during the 2022 bear market, when
Benchmark 2 was getting hammered, Strategy 2 was sitting in cash.
When the market recovered, it came back in. The NAV curve in Figure~\ref{fig:d5}
shows this clearly --- Strategy 2 is the smoothest line in the chart.

\textbf{Which one is better?}
It depends on what you care about.
If you want maximum return and can stomach a drawdown of 30\%, Strategy 1 is the pick.
If protecting capital and sleeping at night matters more, Strategy 2 is clearly superior.

% =============================================================================
\section{Conclusion}
% =============================================================================

This project produced a fully working backtesting engine along with a library of
eleven trading strategies spanning momentum, mean reversion, and hybrid approaches.
The system is clean enough that adding a new strategy takes maybe 30 lines of code,
and everything --- charts, tables, logs --- gets saved automatically.

The most interesting things I learned from running the experiments:

\begin{itemize}[leftmargin=*]
    \item \textbf{Strategy type has to match the market environment.}
          Momentum crushed mean reversion on NVDA because NVDA spent five years
          in a structural uptrend. The same strategies applied to a range-bound
          stock would likely flip the rankings.

    \item \textbf{Signal quality beats weighting sophistication.}
          Picking the right stocks (trailing return top-10) mattered far more than
          how you weighted them. The SMA crossover signal never caught up to
          trailing return, regardless of weighting scheme.

    \item \textbf{A simple regime filter can dramatically reduce drawdown.}
          Strategy 2's biggest contribution is showing that ``sometimes the best
          trade is no trade.'' Going to cash when the market breadth is bad
          cut the maximum drawdown by 70\% relative to the naive momentum benchmark.

\end{itemize}

\end{document}
